\documentclass{amsart}
% For dealing with references we use the comment environment
\usepackage{verbatim}
\newenvironment{reference}{\comment}{\endcomment}
%\newenvironment{reference}{}{}
\newenvironment{slogan}{\comment}{\endcomment}
\newenvironment{history}{\comment}{\endcomment}

% For commutative diagrams we use Xy-pic
\usepackage[all]{xy}

% We use 2cell for 2-commutative diagrams.
\xyoption{2cell}
\UseAllTwocells

% We use multicol for the list of chapters between chapters
\usepackage{multicol}

% This is generally recommended for better output
\usepackage{lmodern}
\usepackage[T1]{fontenc}

% For cross-file-references

% Package for hypertext links:
\usepackage{hyperref}

% For any local file, say "hello.tex" you want to link to please

% Theorem environments.
%
\theoremstyle{plain}
\newtheorem{theorem}[subsection]{Theorem}
\newtheorem{proposition}[subsection]{Proposition}
\newtheorem{lemma}[subsection]{Lemma}

\theoremstyle{definition}
\newtheorem{definition}[subsection]{Definition}
\newtheorem{example}[subsection]{Example}
\newtheorem{exercise}[subsection]{Exercise}
\newtheorem{situation}[subsection]{Situation}

\theoremstyle{remark}
\newtheorem{remark}[subsection]{Remark}
\newtheorem{remarks}[subsection]{Remarks}

\numberwithin{equation}{subsection}

% Macros
%
\def\lim{\mathop{\mathrm{lim}}\nolimits}
\def\colim{\mathop{\mathrm{colim}}\nolimits}
\def\Spec{\mathop{\mathrm{Spec}}}
\def\Hom{\mathop{\mathrm{Hom}}\nolimits}
\def\Ext{\mathop{\mathrm{Ext}}\nolimits}
\def\SheafHom{\mathop{\mathcal{H}\!\mathit{om}}\nolimits}
\def\SheafExt{\mathop{\mathcal{E}\!\mathit{xt}}\nolimits}
\def\Sch{\mathit{Sch}}
\def\Mor{\mathop{\mathrm{Mor}}\nolimits}
\def\Ob{\mathop{\mathrm{Ob}}\nolimits}
\def\Sh{\mathop{\mathit{Sh}}\nolimits}
\def\NL{\mathop{N\!L}\nolimits}
\def\CH{\mathop{\mathrm{CH}}\nolimits}
\def\proetale{{pro\text{-}\acute{e}tale}}
\def\etale{{\acute{e}tale}}
\def\QCoh{\mathit{QCoh}}
\def\Ker{\mathop{\mathrm{Ker}}}
\def\Im{\mathop{\mathrm{Im}}}
\def\Coker{\mathop{\mathrm{Coker}}}
\def\Coim{\mathop{\mathrm{Coim}}}

% Boxtimes
%
\DeclareMathSymbol{\boxtimes}{\mathbin}{AMSa}{"02}

%
% Macros for moduli stacks/spaces
%
\def\QCohstack{\mathcal{QC}\!\mathit{oh}}
\def\Cohstack{\mathcal{C}\!\mathit{oh}}
\def\Spacesstack{\mathcal{S}\!\mathit{paces}}
\def\Quotfunctor{\mathrm{Quot}}
\def\Hilbfunctor{\mathrm{Hilb}}
\def\Curvesstack{\mathcal{C}\!\mathit{urves}}
\def\Polarizedstack{\mathcal{P}\!\mathit{olarized}}
\def\Complexesstack{\mathcal{C}\!\mathit{omplexes}}
% \Pic is the operator that assigns to X its picard group, usage \Pic(X)
% \Picardstack_{X/B} denotes the Picard stack of X over B
% \Picardfunctor_{X/B} denotes the Picard functor of X over B
\def\Pic{\mathop{\mathrm{Pic}}\nolimits}
\def\Picardstack{\mathcal{P}\!\mathit{ic}}
\def\Picardfunctor{\mathrm{Pic}}
\def\Deformationcategory{\mathcal{D}\!\mathit{ef}}

\setlength{\emergencystretch}{3em}
\begin{document}
\title[Stacks Project, Tag 01A9]{710 --- Maps agreeing below degree n between n-coskeletal simplicial sets are homotopic}
\maketitle
\noindent Copyright (C) 2005--2025 Johan de Jong. Stacks Project authors. Source: \href{https://stacks.math.columbia.edu/tag/01A9}{Tag 01A9}.

\noindent Editorial difficulty note (not author text). Difficulty H2: The author gives both a trivial-Kan-fibration lifting construction and an explicit truncated homotopy, with a complete degeneracy and endpoint case analysis. The coskeleton hypothesis provides a specialized extension mechanism beyond qualification topology..

\noindent Editorial provenance note (not author text). History: excerpt from revision \texttt{a0444\allowbreak 6e57e\allowbreak c1fbc\allowbreak 252a8\allowbreak 71afc\allowbreak ec775\allowbreak 2fb28\allowbreak 07b14}, simplicial.tex, lines 6043--6148. Reference presentation was adapted; original-author context and core support blocks were retained or added. The mathematical wording is unchanged. Licensed under GNU FDL 1.2 or later; no invariant sections or cover texts. See ../licenses/COPYING. Context blocks reproduce author definitions and supporting results; they are not additional counted problems.

\begin{definition}
\label{definition-homotopy}
Let $\mathcal{C}$ be a category having finite coproducts.
Suppose that $U$ and $V$ are two simplicial objects of $\mathcal{C}$.
Let $a, b : U \to V$ be two morphisms.
\begin{enumerate}
\item We say a morphism
$$
h : U \times \Delta[1] \longrightarrow V
$$
is a {\it homotopy from $a$ to $b$} if $a = h \circ e_0$ and
$b = h \circ e_1$.
\item We say the morphisms $a$ and $b$ are {\it homotopic} or are
{\it in the same homotopy class}
if there exists a sequence of morphisms $a = a_0, a_1, \ldots, a_n = b$
from $U$ to $V$ such that for each $i = 1, \ldots, n$ there either exists
a homotopy from $a_{i - 1}$ to $a_i$ or there exists a homotopy
from $a_i$ to $a_{i - 1}$.
\end{enumerate}
\end{definition}

\begin{lemma}
\label{lemma-homotopy}
Let $f^0, f^1 : V \to U$ be maps of simplicial sets.
Let $n \geq 0$ be an integer.
Assume
\begin{enumerate}
\item The maps $f^j_i : V_i \to U_i$, $j = 0, 1$ are equal for $i < n$.
\item The canonical morphism $U \to \text{cosk}_n \text{sk}_n U$
is an isomorphism.
\item The canonical morphism $V \to \text{cosk}_n \text{sk}_n V$
is an isomorphism.
\end{enumerate}
Then $f^0$ is homotopic to $f^1$.
\end{lemma}

\begin{proof}[First proof]
Let $W$ be the $n$-truncated simplicial set with $W_i = U_i$ for
$i < n$ and $W_n = U_n / \sim$ where $\sim$ is the equivalence relation
generated by $f^0(y) \sim f^1(y)$ for $y \in V_n$. This makes sense
as the morphisms $U(\varphi) : U_n \to U_i$ corresponding to
$\varphi : [i] \to [n]$ for $i < n$ factor through the quotient map
$U_n \to W_n$ because $f^0$ and $f^1$ are morphisms of simplicial sets and
equal in degrees $< n$. Next, we upgrade $W$ to a simplicial set
by taking $\text{cosk}_n W$. By Lemma \href{https://stacks.math.columbia.edu/tag/01A6}{lemma section (Tag 01A6)}
the morphism $g : U \to W$ is a trivial Kan fibration. Observe
that $g \circ f^0 = g \circ f^1$ by construction and denote
this morphism $f : V \to W$. Consider the diagram
$$
\xymatrix{
\partial \Delta[1] \times V \ar[rr]_{f^0, f^1} \ar[d] & & U \ar[d] \\
\Delta[1] \times V \ar[rr]^f \ar@{-->}[rru] & & W
}
$$
By Lemma \href{https://stacks.math.columbia.edu/tag/08NM}{lemma trivial kan (Tag 08NM)} the dotted arrow exists and the proof is done.
\end{proof}

\begin{proof}[Second proof]
We have to construct a morphism of simplicial sets
$h : V \times \Delta[1] \to U$ which recovers $f^i$ on composing with $e_i$.
The case $n = 0$ was dealt with above the lemma.
Thus we may assume that $n \geq 1$.
The map $\Delta[1] \to \text{cosk}_1 \text{sk}_1 \Delta[1]$
is an isomorphism, see Lemma \href{https://stacks.math.columbia.edu/tag/018E}{lemma simplex cosk (Tag 018E)}.
Thus we see that $\Delta[1] \to \text{cosk}_n \text{sk}_n \Delta[1]$
is an isomorphism as $n \geq 1$, see
Lemma \href{https://stacks.math.columbia.edu/tag/018B}{lemma cosk up (Tag 018B)}. And hence $V \times \Delta[1] \to
\text{cosk}_n \text{sk}_n (V \times \Delta[1])$
is an isomorphism too, see Lemma \href{https://stacks.math.columbia.edu/tag/018C}{lemma cosk product (Tag 018C)}.
In other words, in order to construct the homotopy
it suffices to construct a suitable
morphism of $n$-truncated simplicial sets
$h : \text{sk}_n V \times \text{sk}_n \Delta[1] \to \text{sk}_n U$.

\medskip\noindent
For $k = 0, \ldots, n - 1$ we define $h_k$ by the
formula $h_k(v, \alpha) = f^0(v) = f^1(v)$.
The map $h_n : V_n \times \Mor_{\Delta}([n], [1]) \to U_n$
is defined as follows. Pick $v \in V_n$ and $\alpha : [n] \to [1]$:
\begin{enumerate}
\item If $\Im(\alpha) = \{0\}$, then we set $h_n(v, \alpha) = f^0(v)$.
\item If $\Im(\alpha) = \{0, 1\}$, then we set $h_n(v, \alpha) = f^0(v)$.
\item If $\Im(\alpha) = \{1\}$, then we set $h_n(v, \alpha) = f^1(v)$.
\end{enumerate}
Let $\varphi : [k] \to [l]$ be a morphism of $\Delta_{\leq n}$.
We will show that the diagram
$$
\xymatrix{
V_{l} \times \Mor([l], [1]) \ar[r] \ar[d] & U_{l} \ar[d] \\
V_{k} \times \Mor([k], [1]) \ar[r] & U_{k}
}
$$
commutes.
Pick $v \in V_{l}$ and $\alpha : [l] \to [1]$.
The commutativity means that
$$
h_k(V(\varphi)(v), \alpha \circ \varphi)
=
U(\varphi)(h_l(v, \alpha)).
$$
In almost every case this holds because
$h_k(V(\varphi)(v), \alpha \circ \varphi) = f^0(V(\varphi)(v))$
and $U(\varphi)(h_l(v, \alpha)) = U(\varphi)(f^0(v))$, combined
with the fact that $f^0$ is a morphism of simplicial sets.
The only cases where this does not hold is when
either (A) $\Im(\alpha) = \{1\}$ and $l = n$
or (B) $\Im(\alpha \circ \varphi) = \{1\}$ and $k = n$.
Observe moreover that necessarily $f^0(v) = f^1(v)$
for any degenerate $n$-simplex of $V$.
Thus we can narrow the cases above down even further
to the cases (A) $\Im(\alpha) = \{1\}$, $l = n$
and $v$ nondegenerate, and (B)
$\Im(\alpha \circ \varphi) = \{1\}$, $k = n$
and $V(\varphi)(v)$ nondegenerate.

\medskip\noindent
In case (A), we see that also $\Im(\alpha \circ \varphi) = \{1\}$.
Hence we see that not only $h_l(v, \alpha) = f^1(v)$ but also
$h_k(V(\varphi)(v), \alpha \circ \varphi) = f^1(V(\varphi)(v))$.
Thus we see that the relation holds because $f^1$ is a morphism
of simplicial sets.

\medskip\noindent
In case (B) we conclude that $l = k = n$ and
$\varphi$ is bijective, since otherwise $V(\varphi)(v)$
is degenerate. Thus $\varphi = \text{id}_{[n]}$, which is a trivial case.
\end{proof}
\end{document}
